\section{SpanVouch Reference Implementation}
\label{sec:implementation}

\input{figures/ivad-system}

\subsection{Architecture and Contracts}

\system implements IVAD as a Python system with dependency direction
\code{contracts} $\leftarrow$ \code{trace} $\leftarrow$ \code{diagnosis}
$\leftarrow$ \code{verification} $\leftarrow$ \code{review}. Delivery,
storage, framework, and model integrations remain outside these core layers.
Architecture tests enforce the direction so that a provider or orchestration
framework cannot redefine the evidence and decision semantics.

The public Contract v1 surface has six roots: \traceir, diagnostic context,
diagnosis, verification, review, and artifact manifest. Models reject unknown
fields and use canonical JSON for stable SHA-256 identities. A package release
does not implicitly change a contract version, which allows stored traces,
reports, and decisions to retain their original interpretation.

\subsection{Trace Projection and Evidence Binding}

\code{TraceProjector} converts a raw execution into a sanitized diagnostic
context while preserving stable span identity and causal structure.
\code{EvidenceCatalog} resolves selectors of the form
\code{span-id::\allowbreak{}field-path} to a canonical value and digest. Provider-visible
views exclude credentials, authorization headers, raw provider responses,
hidden reasoning, evaluator labels, mutation metadata, and split identity.

\code{DiagnosisEngine} produces a bounded report containing status, failure
type, decisive spans, at most three staged causal claims, evidence references,
confidence, and provenance. The deterministic rules diagnoser supplies the
offline path; a provider-backed diagnoser uses the same output contract and
must pass strict parsing before its report enters verification.

\subsection{Separated Verification}

\code{DeterministicVerifier} re-resolves every cited selector from the stored
context. It checks report binding, canonical values and hashes, duplicate
references, causal and critical-span grounding, evidence budgets, temporal and
scope constraints, clean-trace conflicts, and domain invariants. These checks
form hard eligibility: a semantic score cannot override a failed integrity or
scope rule.

\code{SemanticVerifier} is an optional provider path that receives only the
sanitized context and structured diagnosis. It evaluates relevance,
sufficiency, counter-evidence, and alternative causal accounts, then returns a
strict verification report. Invalid provider output becomes a typed
human-review requirement rather than triggering an untracked repair call. The
experiment adapters can vary model, provider, prompt, and visible context while
holding the diagnosis bytes fixed.

\subsection{Bounded Review and Recovery}

Verifier reports use three workflow verdicts: \textsc{verified},
\textsc{needs-evidence}, and \textsc{review-required}. The review state machine
permits at most one evidence-guided diagnosis revision and repeats verification
after that revision. Every path terminates in a human \textsc{confirm},
\textsc{correct}, or \textsc{reject} decision; automation does not manufacture
an authenticated reviewer identity.

SQLite persistence stores sanitized snapshots, revisions, verifier reports,
events, and terminal decisions. Idempotency keys, compare-and-swap transitions,
owner-fenced leases, and immutable event order make interrupted cases
recoverable across processes and service restarts. FastAPI and the HTTP-only
command-line client expose the same application workflow. The container image
runs as an unprivileged user and stores review state in a persistent volume.

\subsection{Framework and Evaluation Adapters}

SupportLab and OpsLab create controlled tool-using failures with execution-
derived labels. LangGraph and AutoGen adapters implement a shared runtime port;
parity checks compare terminal state, tool-call order, failure injection, and
TraceIR projection before a trace becomes eligible for paired evaluation. The
replay repository then freezes content-addressed traces so that verifier
conditions consume identical diagnosis and trajectory bytes.

\code{ArtifactManifest} binds code, contracts, datasets, configuration,
prompts, model and provider identities, randomness, runtime, usage, cost, and
payload provenance. The experiment runner fails closed before live execution
unless an exact approved manifest, allowlisted plan, budget ledger, and explicit
live flag agree. Offline rules, deterministic verification, artifact analysis,
and the complete delivery path require no provider credential.

The dashed path in Figure~\ref{fig:ivad-system} marks IVAD's formal
group-selective controller and post-acquisition policy. SpanVouch supplies the
contract and artifact interfaces required by that protocol; the results in this
paper validate the implemented engineering path rather than empirical target-
risk attainment.
